% TOP_PROBLEM: {"id": 3183, "title": "Highly connected graphs with no K_n minor", "category_id": 3, "category": {"id": 3, "name": "graph_theory", "display_name": "Graph Theory", "description": "Problems involving graphs, networks, and their properties.", "slug": "graph-theory", "order_index": 3, "created_at": "2026-07-31T15:26:25.671Z"}, "status": "open", "problem_number": "OPG-153", "set_id": 12, "statusQualification": "Restricts the integer parameter to t>=5 and writes the deletion bound as at most t-5."}
% FIRST_POSTED: 2026-10-01
% ENTRY_KIND: statement_only
% UnsolvedMath ID 3183; source catalogue code OPG-153
% !TeX program = lualatex
% Definition and sources only; this file is not an LLM solution attempt.
\documentclass[11pt,a4paper]{article}
\usepackage[margin=25mm]{geometry}
\usepackage{fontspec}
\setmainfont{lmroman10-regular.otf}[BoldFont=lmroman10-bold.otf,
 ItalicFont=lmroman10-italic.otf,BoldItalicFont=lmroman10-bolditalic.otf]
\usepackage{amsmath,amssymb,mathtools}
\usepackage{unicode-math}
\setmathfont{latinmodern-math.otf}
\AtBeginDocument{\let\setminus\smallsetminus}
\usepackage{xurl,hyperref}
\hypersetup{colorlinks=true,urlcolor=blue}
\setcounter{secnumdepth}{0}
\setlength{\parindent}{0pt}
\setlength{\parskip}{5pt}
\setlength{\emergencystretch}{3em}
\begin{document}
\section*{Highly connected graphs with no K\_n minor}\label{3183}
{\small UnsolvedMath ID 3183 · OPG-153 · Graph Theory · OpenGarden}
\subsection{Definitions and mathematical statement}
For an integer $t\ge5$, a graph is $t$-vertex-connected if it has more than $t$ vertices and deleting any fewer than $t$ vertices leaves it connected. All graphs here are finite and simple. A $K_t$ minor is specified by $t$ disjoint nonempty connected vertex sets, called branch sets, with an edge between every two branch sets. Contracting each branch set and deleting surplus edges yields the complete graph $K_t$.

The conjecture asks whether, for every integer $t\ge5$, there exists an integer $N(t)$ such that every $t$-vertex-connected graph $G$ with $|V(G)|\ge N(t)$ and no $K_t$ minor has a set $S\subseteq V(G)$ satisfying
\[
             |S|\le t-5,\qquad G-S\text{ is planar}.
\]
Planar means that the graph has a drawing in the plane with no crossings between edge interiors. The set $S$ is allowed to depend on $G$, while the order threshold depends only on $t$. The parameter $t$ is the connectivity and excluded-minor parameter, not the order of $G$.

For $t=5$ the deletion set must be empty. For larger $t$, the claim is that sufficiently large examples are obtained from planar graphs by adding at most $t-5$ exceptional vertices, whose incidences need not satisfy any further restriction. The sufficiently-large-order qualification is part of the assertion and must not be replaced by a claim covering every order.
\subsection{Short English statement}
For each fixed connectivity t, are all sufficiently large t-connected graphs without a complete t-vertex minor planar after deleting at most t−5 vertices?
\subsection{Sources}
\begin{itemize}
\item[S1] ULAM.AI and the UnsolvedMath Contributors, supplied problems.json, record 3183 (OPG-153), OpenGarden. Catalogue identity and source transcription; CC BY 4.0.
\url{https://huggingface.co/datasets/ulamai/UnsolvedMath}.
\item[S2] Open Problem Garden, Highly connected graphs with no K\_n minor. Original problem and discussion.
\url{http://www.openproblemgarden.org/op/high_connectivity_no_k_n}.
\end{itemize}
\end{document}
